\documentclass[a4paper,11pt]{article}

\usepackage[T1]{fontenc}
\usepackage[utf8]{inputenc}
\usepackage{lmodern}
\usepackage{microtype}
\usepackage[a4paper,margin=1in,headheight=15pt,headsep=14pt]{geometry}
\usepackage{amsmath,amssymb}
\usepackage{graphicx}
\usepackage{booktabs}
\usepackage[table,svgnames]{xcolor}
\usepackage[most]{tcolorbox}
\tcbuselibrary{skins,breakable}
\usepackage{pifont}

\IfFileExists{bbding.sty}{%
  \usepackage{bbding}%
  \newcommand{\glyphHand}{\Pointinghand}%
  \newcommand{\glyphPencil}{\Pencil}%
}{%
  \newcommand{\glyphHand}{\ding{43}}%
  \newcommand{\glyphPencil}{\ding{46}}%
}

\usepackage{enumitem}
\setlist{leftmargin=1.5em,topsep=4pt,itemsep=3pt}

\newlist{steps}{enumerate}{1}
\setlist[steps]{label=\textbf{Step \arabic*.},
  leftmargin=4.4em, labelwidth=3.8em, labelsep=0.6em, labelindent=0pt,
  align=left, topsep=5pt, itemsep=6pt}
\usepackage{parskip}
\usepackage{fancyhdr}
\usepackage{titlesec}
\usepackage[hidelinks]{hyperref}
\usepackage{xurl}
\hypersetup{breaklinks=true}

\definecolor{navy}{HTML}{21355E}
\definecolor{navyrule}{HTML}{21355E}
\definecolor{inktext}{HTML}{1A202C}
\definecolor{mutedtext}{HTML}{6B7280}

\definecolor{intuFrame}{HTML}{2C5AA0}   \definecolor{intuFill}{HTML}{F4F7FC}
\definecolor{everyFrame}{HTML}{8A5A1E}  \definecolor{everyFill}{HTML}{FBF1E3}
\definecolor{workFrame}{HTML}{2E7D52}   \definecolor{workFill}{HTML}{F1F8F3}
\definecolor{watchFrame}{HTML}{B23A48}  \definecolor{watchFill}{HTML}{FBF1F2}
\definecolor{keyFrame}{HTML}{6A4C93}    \definecolor{keyFill}{HTML}{F4F0F9}
\definecolor{waitFrame}{HTML}{0F766E}   \definecolor{waitFill}{HTML}{EEF7F5}

\color{inktext}
\hypersetup{linkcolor=navy,urlcolor=navy,citecolor=navy}

\newcommand{\CourseShort}{Machine Learning}
\newcommand{\SessionNo}{12}
\newcommand{\LectureTitle}{Ensemble Methods \& Text Classification}

\pagestyle{fancy}
\fancyhf{}
\fancyhead[L]{\footnotesize\color{mutedtext}\textsf{\CourseShort}}
\fancyhead[R]{\footnotesize\color{mutedtext}\textsf{Session \SessionNo\ \textperiodcentered\ \LectureTitle}}
\fancyfoot[L]{\footnotesize\color{mutedtext}\textsf{WILP, BITS Pilani}}
\fancyfoot[C]{\footnotesize\color{mutedtext}\thepage}
\renewcommand{\headrulewidth}{0.4pt}
\renewcommand{\footrulewidth}{0pt}
\renewcommand{\headrule}{\hbox to\headwidth{\color{navyrule!55}\leaders\hrule height \headrulewidth\hfill}}

\titleformat{\section}
  {\sffamily\Large\bfseries\color{navy}}
  {\thesection}
  {0.8em}
  {}
  [{\vspace{2pt}\color{navy!70}\titlerule[0.8pt]}]
\titleformat{\subsection}
  {\sffamily\large\bfseries\color{navy!92}}
  {\thesubsection}{0.7em}{}
\titlespacing*{\section}{0pt}{16pt}{8pt}
\titlespacing*{\subsection}{0pt}{12pt}{4pt}

\newif\ifmlkfirstsec \mlkfirstsectrue
\newcommand{\sectionbreak}{\ifmlkfirstsec\mlkfirstsecfalse\else\clearpage\fi}

\newcommand{\secsub}[1]{%
  \noindent{\sffamily\bfseries\itshape\color{navy!85}\large #1}\par\medskip}

\newcommand{\flipanswer}[1]{\rotatebox[origin=c]{180}{\parbox{0.92\linewidth}{\small #1}}}

\tcbset{
  housecallout/.style 2 args={
    enhanced, breakable,
    colback=#2, colframe=#1,
    boxrule=0.6pt, arc=2.4pt, outer arc=2.4pt,
    left=10pt, right=10pt, top=9pt, bottom=9pt,
    before skip=11pt, after skip=11pt,
    fontupper=\normalsize,
    coltitle=white,
    fonttitle=\bfseries\sffamily\small,
    attach boxed title to top left={xshift=6mm,yshift=-3mm},
    boxed title style={
      colback=#1, colframe=#1,
      rounded corners, arc=3pt,
      boxrule=0pt, left=6pt, right=6pt, top=2pt, bottom=2pt
    }
  }
}

\newtcolorbox{intuition}{housecallout={intuFrame}{intuFill}, title={\raisebox{-0.05ex}{\glyphHand}\,\ The intuition}}
\newtcolorbox{everyday}{housecallout={everyFrame}{everyFill}, title={\raisebox{-0.05ex}{$\bigstar$}\,\ Everyday picture}}
\newtcolorbox{worked}[1]{housecallout={workFrame}{workFill}, title={\raisebox{-0.05ex}{\glyphPencil}\,\ Worked example: #1}}
\newtcolorbox{watchout}{housecallout={watchFrame}{watchFill}, title={\raisebox{-0.05ex}{\ding{55}}\,\ Watch out}}
\newtcolorbox{keytake}{housecallout={keyFrame}{keyFill}, title={\raisebox{-0.05ex}{\ding{51}}\,\ Key takeaway}}
\newtcolorbox{waitwhat}{housecallout={waitFrame}{waitFill}, title={\raisebox{-0.05ex}{\ding{93}}\,\ Wait --- really?}}

\newcommand{\bitswordmark}{{\sffamily\bfseries\color{white}\fontsize{12.5}{15}\selectfont WILP, BITS Pilani}}
\newcommand{\bitslogochip}{}

\newcommand{\titlebanner}[4]{%
  \begin{tcolorbox}[
    enhanced, breakable=false,
    colback=navy, colframe=navy,
    boxrule=0pt, arc=3pt, outer arc=3pt,
    left=16pt, right=16pt, top=16pt, bottom=16pt,
    before skip=2pt, after skip=14pt,
    width=\linewidth ]
    \noindent
    \begin{minipage}[c]{0.72\linewidth}\bitswordmark\end{minipage}%
    \hfill
    \begin{minipage}[c]{0.26\linewidth}\raggedleft\bitslogochip\end{minipage}\par
    \vspace{9pt}
    {\sffamily\footnotesize\bfseries\color{white!85}\MakeUppercase{#1}}\par
    \vspace{6pt}
    {\sffamily\bfseries\fontsize{26}{30}\selectfont\color{white} #2\par}
    \vspace{5pt}
    {\sffamily\large\color{white!88} #3\par}
    \vspace{9pt}
    {\color{white!55}\rule{\linewidth}{0.4pt}}\par
    \vspace{6pt}
    {\sffamily\footnotesize\color{white!82} #4\par}
  \end{tcolorbox}%
}

\newcommand{\housefig}[2]{%
  \begin{figure}[htbp]
    \centering
    \includegraphics[width=\linewidth]{#1}%
    \caption{#2}
  \end{figure}%
}

\newcommand{\vocab}[1]{\textbf{\color{navy}#1}}

\begin{document}

\thispagestyle{fancy}

\titlebanner
  {Machine Learning \textperiodcentered\ Companion Reader}
  {Ensemble Methods \& Text Classification}
  {Committee Models, Bagging, Boosting, Multinomial Naïve Bayes, and the Connection to Logistic Regression}
  {Session 12 \textperiodcentered\ Read alongside the lecture slides \textperiodcentered\ Every slide example worked out in full}

\smallskip
\noindent{\small\color{mutedtext}Prepared for the Work Integrated Learning Programmes (WILP), BITS Pilani.}
\smallskip

\section{Ensemble Philosophy \& Committee Models}
\secsub{Combining multiple diverse base classifiers reduces variance and produces superior predictive performance.}

According to the \vocab{No Free Lunch Theorem}, no single learning algorithm outperforms all others across all tasks. \vocab{Ensemble Learning} constructs a committee of base models $C_1, C_2, \dots, C_T$ and aggregates their predictions.

\begin{intuition}
Wisdom of the crowd: A committee of diverse, independent decision-makers makes far fewer errors than any single individual expert.
\end{intuition}

\housefig{figures/fig_ensemble_variance.png}{Bagging ensemble averaging reduces variance while maintaining low model bias.}

\section{Bagging (Bootstrap Aggregating)}
\secsub{Bagging creates diverse base models by training each on an independent bootstrap sample of the dataset.}

Given a training set $D$ of size $N$:
\begin{enumerate}
  \item Generate $T$ bootstrap samples $D_1, D_2, \dots, D_T$ by sampling $N$ instances from $D$ \vocab{with replacement}.
  \item On average, each bootstrap sample contains $1 - (1 - 1/N)^N \approx 63.2\%$ of unique original instances, leaving $36.8\%$ as Out-Of-Bag (OOB) validation data.
  \item Train base classifier $C_t$ on $D_t$.
  \item Predict using majority voting (classification) or averaging (regression).
\end{enumerate}

\begin{worked}{1D Bagging Ensemble Voting Across 8 Rounds}
Consider 10 points $x \in \{0.1, 0.2, 0.3, 0.4, 0.5, 0.6, 0.7, 0.8, 0.9, 1.0\}$.
Decision stump base classifiers pick split thresholds across 8 bagging rounds:
\begin{steps}
  \item \textbf{Round 1 (Split $x=0.35$):} Predicts $+1$ for $x \le 0.35$, $-1$ for $x > 0.35$.
  \item \textbf{Round 2 (Split $x=0.70$):} Predicts $+1$ for $x \le 0.70$, $+1$ for $x > 0.70$.
  \item \textbf{Round 6 (Split $x=0.75$):} Predicts $-1$ for $x \le 0.75$, $+1$ for $x > 0.75$.
  \item \textbf{Majority Vote Aggregation:} Taking the mode vote across all 8 rounds smooths out local variance and produces a stable, solid decision boundary.
\end{steps}
\end{worked}

\section{Multinomial Naïve Bayes for Text Classification}
\secsub{Multinomial Naïve Bayes models document term frequencies for text categorization.}

In the \vocab{Bag of Words (BoW)} model, text order is ignored and word counts $f_i$ represent the document. The posterior class score is:
\begin{equation}
P(c \mid d) \propto P(c) \prod_{i=1}^{|V|} P(w_i \mid c)^{f_i}, \quad P(w_i \mid c) = \frac{N_{c, w_i} + 1}{N_c + |V|}.
\end{equation}

\housefig{figures/fig_multinomial_nb.png}{Multinomial word likelihood probabilities P(w|c) for Sports vs Not Sports text classification.}

\begin{worked}{Sentence Tagging: "A very close game"}
Classify sentence $d = \text{"A very close game"}$ into $c \in \{\text{Sports}, \text{Not Sports}\}$.
Vocabulary size $|V| = 6$. Total words in Sports $N_{\text{Sports}} = 8$, in Not Sports $N_{\text{NotSports}} = 3$.
Word count in Sports: $N_{\text{Sports, close}} = 1$, $N_{\text{Sports, game}} = 1$.
\begin{steps}
  \item \textbf{Compute Smoothed Word Probabilities for Sports:}
        \begin{equation}
        P(\text{close} \mid \text{Sports}) = \frac{1 + 1}{8 + 6} = \frac{2}{14} = \frac{1}{7}, \quad P(\text{game} \mid \text{Sports}) = \frac{1 + 1}{8 + 6} = \frac{2}{14} = \frac{1}{7}.
        \end{equation}
  \item \textbf{Compute Smoothed Word Probabilities for Not Sports:}
        \begin{equation}
        P(\text{close} \mid \text{NotSports}) = \frac{0 + 1}{3 + 6} = \frac{1}{9}, \quad P(\text{game} \mid \text{NotSports}) = \frac{0 + 1}{3 + 6} = \frac{1}{9}.
        \end{equation}
  \item \textbf{Compare Class Scores:}
        \begin{align}
        P(\text{Sports} \mid d) &\propto 0.5 \times \frac{1}{7} \times \frac{1}{7} = \frac{1}{98} \approx 0.0204,\\
        P(\text{NotSports} \mid d) &\propto 0.5 \times \frac{1}{9} \times \frac{1}{9} = \frac{1}{162} \approx 0.0123.
        \end{align}
        Since $0.0204 > 0.0123$, the sentence is correctly classified as \textbf{Sports}!
\end{steps}
\end{worked}

\section{Mathematical Link: Gaussian Naïve Bayes \& Logistic Regression}
\secsub{Gaussian Naïve Bayes with equal variances mathematically implies the Logistic Regression functional form.}

For binary classification $Y \in \{0, 1\}$, assume features $X_i \mid Y=y \sim \mathcal{N}(\mu_{iy}, \sigma_i^2)$ with equal class variances $\sigma_{i0} = \sigma_{i1} = \sigma_i$.

\housefig{figures/fig_gnb_logistic_link.png}{Logistic sigmoid curve derived from Gaussian Naïve Bayes log-odds ratio.}

Applying Bayes' rule:
\begin{align}
P(Y=1 \mid \mathbf{X}) &= \frac{P(Y=1) P(\mathbf{X} \mid Y=1)}{P(Y=1) P(\mathbf{X} \mid Y=1) + P(Y=0) P(\mathbf{X} \mid Y=0)} \nonumber\\
&= \frac{1}{1 + \exp\left( -\ln \frac{P(Y=1)P(\mathbf{X}\mid Y=1)}{P(Y=0)P(\mathbf{X}\mid Y=0)} \right)}.
\end{align}

Expanding the log-odds ratio with Gaussian densities:
\begin{equation}
\ln \frac{P(Y=1)}{P(Y=0)} + \sum_{i=1}^d \ln \frac{p(X_i \mid Y=1)}{p(X_i \mid Y=0)} = \theta_0 + \sum_{i=1}^d \theta_i X_i = \boldsymbol{\theta}^T \mathbf{X},
\end{equation}
where $\theta_i = \frac{\mu_{i1} - \mu_{i0}}{\sigma_i^2}$. So:
\begin{equation}
P(Y=1 \mid \mathbf{X}) = \frac{1}{1 + e^{-\boldsymbol{\theta}^T \mathbf{X}}}.
\end{equation}
Gaussian Naïve Bayes directly implies the exact sigmoid functional form of Logistic Regression!

\section*{Glossary \& Symbols}
\addcontentsline{toc}{section}{Glossary \& Symbols}

\noindent\textbf{Key Terms}
\begin{description}[leftmargin=8.5em,style=nextline,font=\normalfont\bfseries\color{navy}]
  \item[Ensemble Learning] Combining predictions from multiple base models to improve stability and accuracy.
  \item[Bagging] Bootstrap Aggregating; sampling data with replacement to reduce variance.
  \item[Multinomial Naïve Bayes] Naïve Bayes model for discrete word frequency distributions.
  \item[Out-Of-Bag (OOB)] The $\approx 36.8\%$ of training samples excluded from a bootstrap sample.
\end{description}

\medskip
\noindent\textbf{Symbols at a Glance}
\begin{center}\small
\begin{tabular}{@{}ll@{}}
\toprule
\textbf{Symbol} & \textbf{Meaning} \\
\midrule
$C_t$ & Base classifier model $t$ in ensemble \\
$|V|$ & Vocabulary size in text classification \\
$\boldsymbol{\theta}^T \mathbf{X}$ & Linear log-odds decision boundary score \\
\bottomrule
\end{tabular}
\end{center}

\end{document}
